% 由 scripts/make_paper_tables.py 从 results/ct_all/*.json 生成，请勿手改
\begin{table}[t]
  \centering
  \caption{Output execution versus semantic correctness. The clean residual stream is
           restored at one layer of the corrupted run; we record whether the model emits a
           parseable number (\emph{parse}) and whether that number is correct (\emph{acc}).
           Parseability is essentially layer-invariant, whereas correctness moves with the
           restored depth --- and only for problems the model can solve without corruption.}
  \label{tab:layer-parse-acc}
  \small
  \begin{tabular}{@{}l c c c@{}}
    \toprule
    \multirow{2}{*}{\textbf{Condition}} & \multirow{2}{*}{\textbf{Parse rate}} &
      \multicolumn{2}{c}{\textbf{Accuracy}} \\
    \cmidrule(lr){3-4}
    & & \textbf{clean correct (n=10)} & \textbf{clean wrong (n=90)} \\
    \midrule
    corrupted (no restore) & 96.0\% & -- & 1.0\% \\
    clean (no corruption)  & 100.0\% & -- & 10.0\% \\
    \midrule
    restore layer 0 & 99.0\% & 70.0\% & 1.1\% \\
    restore layer 4 & 99.0\% & 40.0\% & 3.3\% \\
    restore layer 8 & 99.0\% & 50.0\% & 0.0\% \\
    restore layer 12 & 97.0\% & 70.0\% & 0.0\% \\
    restore layer 16 & 100.0\% & 60.0\% & 3.3\% \\
    restore layer 20 & 100.0\% & 30.0\% & 2.2\% \\
    restore layer 24 & 97.0\% & 10.0\% & 0.0\% \\
    \bottomrule
  \end{tabular}
  \vspace{2pt}
  \parbox{\linewidth}{\footnotesize
  Qwen2.5-7B-Instruct, 100 attempted problems, 7 restore points
  (every 4 layers of 28). Parse rate stays within
  97.0--100.0\% at every layer, including the corrupted baseline:
  the model always emits a number. Correctness, by contrast, is only recoverable for
  problems the clean run solves, and never recovers for the rest at any single layer.}
\end{table}
